\DocumentMetadata{
  pdfversion=2.0,
  pdfstandard=ua-2,
  lang=en-US,
  tagging=on,
  tagging-setup={math/setup=mathml-SE,tabsorder=structure}
}
\documentclass[11pt,letterpaper]{article}
\usepackage[margin=0.68in,includefoot]{geometry}
\usepackage{newcomputermodern}
\usepackage{amsmath,amscd,mathtools}
\usepackage{tikz,adjustbox}
\providecommand{\square}{\mdlgwhtsquare}
\usepackage[hidelinks]{hyperref}
\hypersetup{
  pdftitle={Topology First-Year Exam, May 2023, Part 1},
  pdfauthor={University of Florida Department of Mathematics},
  pdfsubject={Graduate mathematics examination},
  pdfdisplaydoctitle=true,
  bookmarksopen=true
}
\setcounter{secnumdepth}{0}
\setlength{\parindent}{0pt}
\setlength{\parskip}{0.28em}
\setlength{\emergencystretch}{3em}
\setlength{\leftmargini}{2.15em}
\setlength{\leftmarginii}{2.65em}
\setlength{\leftmarginiii}{3em}
\pagestyle{empty}
\raggedbottom
\allowdisplaybreaks
\NewTaggingSocketPlug{block/list/label}{examproblem}{%
  \tagstructbegin{tag=Lbl}\tagstructend%
  \tagstructbegin{tag=\LBody}%
  \tagstructbegin{tag=\ProblemHeadingTag,title-o={\ProblemHeadingText}}%
  \tagmcbegin{tag=\ProblemHeadingTag,actualtext={\ProblemHeadingText}}%
  #2%
  \tagmcend%
  \tagstructend%
}
\newenvironment{examproblems}[2]{%
  \def\ProblemHeadingTag{#2}%
  \AssignTaggingSocketPlug{block/list/label}{examproblem}%
  \begin{enumerate}%
  \setcounter{enumi}{#1}%
  \renewcommand{\labelenumi}{\textbf{\theenumi.}}%
  \setlength{\topsep}{0.35em}%
  \setlength{\partopsep}{0pt}%
  \setlength{\itemsep}{0.48em}%
  \setlength{\parsep}{0.18em}%
}{\end{enumerate}}
\newenvironment{examsubparts}{%
  \AssignTaggingSocketPlug{block/list/label}{default}%
  \begin{enumerate}%
  \setlength{\topsep}{0.18em}%
  \setlength{\partopsep}{0pt}%
  \setlength{\itemsep}{0.16em}%
  \setlength{\parsep}{0.1em}%
}{\end{enumerate}}
\newenvironment{examinstructions}{%
  \par\begingroup\itshape
}{%
  \par\endgroup\vspace{0.18em}
}
\makeatletter
\renewcommand\subsection{\@startsection{subsection}{2}{0pt}%
  {-1.25ex plus -.25ex minus -.1ex}%
  {0.55ex plus .1ex}%
  {\normalfont\large\bfseries}}
\renewcommand{\@maketitle}{%
  \newpage\null\vskip -1em%
  {\footnotesize\bfseries
    \href{https://www.ufl.edu/}{University of Florida}, \href{https://math.ufl.edu/}{Department of Mathematics}\par}%
  \vskip -0.5em%
  \tagstructbegin{tag=H1,title={Topology First-Year Exam, May 2023, Part 1}}%
  \tagpdfparaOff%
  \tagmcbegin{tag=H1}%
    {\LARGE\bfseries\href{https://gma.math.ufl.edu/exams/topology-first-year/}{Topology First-Year Exam}, May 2023, Part 1\par}%
  \tagmcend%
  \tagpdfparaOn%
  \tagstructend%
  \vskip 0.25em%
  \tagmcbegin{artifact}%
    \hrule height 1pt%
  \tagmcend%
  \vskip 1em%
}
\makeatother
\title{Topology First-Year Exam, May 2023, Part 1}
\author{}
\date{}
\begin{document}
\maketitle
\thispagestyle{empty}
\begin{examinstructions}
Work the following problems and show all work. Support all statements to the best of your ability.
\end{examinstructions}
\begin{examproblems}{0}{H2}
\def\ProblemHeadingText{Problem 1.}
\item
Show that every continuous map \(f:S^1\to S^2\) of the circle to the 2-sphere is null-homotopic.
\def\ProblemHeadingText{Problem 2.}
\item
Show that the Stone–Čech compactification \(\beta(\mathbb{N})\) of the naturals is uncountable.
\def\ProblemHeadingText{Problem 3.}
\item
Let \(X_{\geq 0}\), \(Y_{\geq 0}\), and \(Z_{\geq 0}\) denote the nonnegative parts of the~\(x\)-axis, the~\(y\)-axis, and the~\(z\)-axis in \(\mathbb{R}^3\) respectively. Let \(A=\mathbb{R}^3\setminus(X_{\geq 0}\cup Y_{\geq 0}\cup Z_{\geq 0})\). Compute the fundamental group \(\pi_1(A)\).
\def\ProblemHeadingText{Problem 4.}
\item
Let \(f:S^1\to T=S^1\times S^1\) be a map to a torus defined by the formula \(f(z)=(z,-z)\). Compute the fundamental group \(\pi_1(X)\) where \(X=T\cup_f B^2\) is obtained from~\(T\) by attaching a 2-disk along~\(f\).
\def\ProblemHeadingText{Problem 5.}
\item
Show that the projective plane \(\mathbb{R}P^2\) is not homeomorphic to the Klein bottle~\(K\). What about the connected sum \(\mathbb{R}P^2\#\mathbb{R}P^2\)?
\end{examproblems}
\begin{examinstructions}
Answer the following with complete definitions or statements or short proofs.
\end{examinstructions}
\begin{examproblems}{5}{H2}
\def\ProblemHeadingText{Problem 6.}
\item
Show that \(S^2\) is not homeomorphic to \(S^3\).
\def\ProblemHeadingText{Problem 7.}
\item
Let \(CX\) denote the cone over the~\(n\)-point space \(X=\{1,\ldots,n\}\). Show that every continuous map \(f:CX\to CX\) has a fixed point.
\def\ProblemHeadingText{Problem 8.}
\item
State the Borsuk–Ulam Theorem for \(S^2\).
\def\ProblemHeadingText{Problem 9.}
\item
Show that if \(g:S^2\to S^2\) is continuous and \(g(x)\neq g(-x)\) for all~\(x\), then~\(g\) is surjective.
\def\ProblemHeadingText{Problem 10.}
\item
State the Tychonoff theorem.
\def\ProblemHeadingText{Problem 11.}
\item
State the Urysohn Lemma.
\def\ProblemHeadingText{Problem 12.}
\item
State the Seifert–van Kampen Theorem. Can it be used to compute the fundamental group of \(S^1\)?
\def\ProblemHeadingText{Problem 13.}
\item
Show that for a retraction \(r:X\to A\), for any \(a_0\in A\), the induced map \(r_*:\pi_1(X,a_0)\to\pi_1(A,a_0)\) on the fundamental groups is surjective.
\def\ProblemHeadingText{Problem 14.}
\item
Give definition of a deformation retraction. Let~\(A\) be a deformation retract of~\(X\). Show that the inclusion map \(A\to X\) is a homotopy equivalence.
\def\ProblemHeadingText{Problem 15.}
\item
Is \(I^I\) separable?
\end{examproblems}
\end{document}
